% Fichier généré automatiquement par tools/generate_corrections.py.
% Source : DYN/DYN-04-TorseurDynamique/STOCK/12_BielleManivelle/12_BielleManivelle.tex
% Statut : complet (2/2 question(s) renseignée(s), 2 reprise(s)).
\begin{corrigebox}[Corrigé extrait de la banque]
\CorrigeQuestion{1}
\CorrigeSource{Réponse reprise d'un exercice homonyme de la banque ; la provenance exacte est indiquée dans le commentaire du fichier.}
% Réponse reprise de : CIN/CIN-02-VitesseAcceleration/12_BielleManivelle/12_BielleManivelle.tex
On commence par calculer $\vectv{B}{2}{0}$ = $\vectv{B}{2}{1}+\vectv{B}{1}{0}$ $=\vectv{B}{1}{0}$.
\begin{itemize}
\item \textbf{Méthode 1 -- dérivation vectorielle :}  $\vectv{B}{1}{0}$ $=\deriv{AB}{\rep{0}}$ 
$=\deriv{R\vi{1}}{\rep{0}}$ 
$= R\thetap(t) \vj{1}$. 
\item \textbf{Méthode 2 -- formule de changement de point: }
$\babarv{B}{A}{1}{0}$ 
$= -R\vi{1}\wedge \thetap{t} \vk{0}$
$= R\thetap(t) \vj{1}$.
\end{itemize}

On a alors, 
$\torseurcin{V}{2}{0}$
$=\torseurl{\varphip(t)\vk{0}}{R\thetap(t) \vj{1}}{B}$.
\CorrigeQuestion{2}
\CorrigeSource{Réponse reprise d'un exercice homonyme de la banque ; la provenance exacte est indiquée dans le commentaire du fichier.}
% Réponse reprise de : CIN/CIN-02-VitesseAcceleration/12_BielleManivelle/12_BielleManivelle.tex
$\vectg{C}{2}{0}$
$=\deriv{\vectv{C}{2}{0}}{\rep{0}}$
$=\lambdapp(t) \vj{0} $.
\end{corrigebox}
